\section{Related Work}
\label{sec:related_work}

\paragraph{Ambiguity in NLP and requirements engineering.}
\citet{li_ambiguity_taxonomy_2024} propose an 11-type NLP taxonomy covering syntactic, semantic, and pragmatic ambiguity, from which we draw our linguistic-layer types; \citet{kamath_scope_2024} find systematic surface-scope preferences in LLMs. In requirements engineering, \citet{massey_ambiguity_re_2014} classify ambiguity types in regulatory texts, and \citet{berry_kamsties_2004} distinguish nocuous ambiguity (stakeholders diverge) from innocuous (they converge). Computational detection has progressed from heuristic \citep{chantree_2006} through NLP-based \citep{ferrari_ambiguity_2019} to transformer-based methods \citep{ezzini_taphsir_2022}. We extend Berry and Kamsties' framework to LLM agents, where the absence of shared organizational context makes previously innocuous ambiguity nocuous.

\paragraph{Psycholinguistics of ambiguity processing.}
\citet{ferreira_good_enough_2002} show that humans build ``good-enough'' representations rather than full parses, paralleling our observation that LLM agents adopt surface interpretations; \citet{li_garden_path_llm_2024} demonstrate analogous garden-path effects in LLMs. \citet{liu_ambiguity_2023} show that language models collapse multiple interpretations into a single output; \citet{gur_arieh_ambiguity_collapse_2026} taxonomize the epistemic risks of this collapse. These findings ground our observation that agents fail to recognize ambiguity---a phenomenon we call \emph{disambiguation neglect} (\S\ref{sec:dn}): they process ambiguous input without meta-cognitive awareness and, unlike humans, lack social channels to recover.

\paragraph{Agent safety and policy specification.}
Agent safety benchmarks evaluate risk awareness \citep{rjudge_2024}, tool-use safety \citep{toolemu_2024}, and policy compliance \citep{yin_safeagentbench_2025,phantompolicy_2026}; our work differs by focusing on the \emph{source specification} rather than agent execution. DRIP-R \citep{drip_r_2026} is the closest concurrent work, benchmarking 400 policy-ambiguous retail scenarios across five RE ambiguity types. Our matched-pair design establishes causality (vs.\ DRIP-R's correlational design), measures safety violations (vs.\ decision quality), and adds linguistic-layer types (Appendix~\ref{app:drip_mapping}); our perturbation controls ($\Delta_{\text{perturb}} = -0.67$\,pp) isolate the causal contribution. SOPBench \citep{li_sopbench_2025} benchmarks SOP compliance with naturally occurring ambiguity; we use it for cross-benchmark evaluation (\S\ref{sec:results}). $\tau$-bench \citep{tau_bench_2024} and $\tau^2$-bench \citep{tau2_bench_2025} provide our evaluation platform. PolicyBank \citep{policybank_2026} studies iterative policy refinement; \citet{krakovna2020specification} document specification gaming in RL agents, a reward-specification analogue of our findings.
